\section{推理任务设计与示例}
\subsection{多步骤推理示例}
要让 self-rewarding pipeline 学会多步骤推理，Weston 建议把 Chain-of-Thought 拆成 `decompose → verify → finalize' 三步。例如：
\begin{itemize}
    \item Prompt：``请分析一个 4 步的代数题，用中间变量展示每一步''。
    \item Draft：列出每分钟的推理链，并给出初始答案。
    \item Verification：让模型重新审阅 draft，并在无法验证时生成 alternative 思路。
\end{itemize}
\textit{``You can make it kind of cross check itself and see that it was wrong and write a final verified response.''}（SRT 1188-1190）这个过程就是 multi-step reasoning + meta judge 的典型应用。

\begin{table}[H]
\centering
\begin{tabular}{p{4cm}p{4cm}p{6cm}}
\toprule
任务类型 & 样例 Prompt & 验证指标 \\
\midrule
代数推理 & ``分解 x+2y=7, x-y=1'' & Chain-of-Thought 完整性、最终答案一致性 \\
因果分析 & ``描述一个 RL agent 如何反思 reward hacking'' & 是否指出操作步骤中的偏差点 \\
统计判断 & ``比较两个模型的精度与置信度'' & Meta judge 判定 winner 的一致性 \\
\bottomrule
\end{tabular}
\caption{Self-Instruct 中典型的 multi-step reasoning 任务。}
\end{table}

\begin{knowledgebox}{多步骤推理的质量控制}
每一条推理链都要配备 verified response，这样 reward 模型能把 errors 视为 supervise signal，而不是只奖励长度或流畅度。过短的 reasoning chain、overthinking 的 chain 都要通过 judge cross-check 扣分。
\end{knowledgebox}

\subsection{编程与验证示例}
在编程任务中，Self-Instruct scene 通常要求模型给出完整代码 + 说明 + 测试。Self-rewarding pipeline 会让模型在生成代码后，再次 ``run'' 自己的推理（例如模拟测试用例），再给出 verified response 说明为什么这段代码满足要求。这样的 verified response 顺带教授模型如何 self-evaluate output。

\begin{itemize}
    \item Prompt：``写一个 Python 函数，输入列表返回唯一元素''。
    \item Draft：输出函数 + 简短解释。
    \item Meta judge：用 SelfT evaluator 比较 draft 与 verified，判定哪份解释更清晰，哪份代码更安全。
\end{itemize}

\subsection{Meta judge 模拟场景}
Meta judge 在实际 deployment 里扮演 ``最后一道 guardrail'' 的角色。例如，当 draft 与 verified 都完成后，judge 会计算二者的 alignment 差值，并给出高可信度的 tick。这种 check 有时需要引入 external agent（如 human or third-party evaluator），特别是在 judge 自信但 hallucinates 的情形（SRT 2500-2509）。

\subsection{本章小结}
通过 multi-step reasoning、编程任务和 meta judge 模拟，可以把 self-rewarding 训练当作一个顶层设计，逐步把 prompt → draft → verified → judge 串联成一个可重复的训练模板。

